\subsection{Problem Statement}

\subsubsection{Key Terminology} \label{sec:key-terms}

\textit{Failures} represent the actual inability of a service to execute its functions~\cite{Soldani2022rcasurvey}. \textit{Faults} correspond to the root causes of such failures (e.g., CPU hog, memory leak, or network disconnection) \cite{Soldani2022rcasurvey, avizienis2004basic}. \textit{Anomalies} are defined as observable symptoms of failures~\cite{Li2022Circa, Soldani2022rcasurvey}. \textit{Root cause analysis (RCA)} is the process of determining why a failure has occurred~\cite{lee2023eadro}, i.e., finding the root cause of the failure. RCA involves a thorough examination of various monitoring data, i.e., including metrics data. \textit{Metrics} are recorded by the monitoring system and contain various critical information within the microservice systems, such as workload, resource consumption, and response time \cite{Xin2023CausalRCA}. These metrics are typically represented as multivariate time series, with each time series corresponding to the data collected with a specific metric. In the context of metric-based RCA, \textit{root cause metrics} are the metrics that are indicators of the root cause \cite{Li2022Circa, chen2022adaptive, Azam2022rcd}. The system operators can use these suggested root cause metrics to identify the true underlying root cause of the failures. The use of these terms aligns with existing RCA works \cite{Chen2014Causeinfer, chen2016causeinfer, chen2022adaptive, Azam2022rcd, Li2022Circa, liu2023pyrca, Meng2020Microcause, Xin2023CausalRCA}.


\subsubsection{Problem Formulation} \label{sec:problem-formulation}
Let us consider a microservice system $\mathcal{S}$ consisting of $n$ services $\{ s^i \}_{i=1}^n$. At each time step $t$, the monitoring system collects $m$ metrics $\mathcal{M}^{i,j=1:m}_t = \{x^{(i, j)}_t\}_{j=1}^{m} (m \geq 1)$ from each service $s^i$. Given a $T$-length observation window with time-series metrics data $\mathcal{M}^{i=1:n,j=1:m}_{t_0:t_0+T}$, our goal is to develop a framework that solves two problems. The first problem is to \textbf{\textit{predict the existence of anomalies}} (failures), represented as a binary indicator $y$, which takes the value of $0$ when there is no anomaly and $1$ when there is an anomaly within the metrics dataset. The second problem is that, when $y$ returns $1$, an RCA module is triggered to \textit{\textbf{pinpoint the root causes of the failure}} using this dataset $\mathcal{M}^{i=1:n,j=1:m}_{t_0:t_0+T}$, i.e., the specific root cause services and the corresponding root cause metrics.

\subsection{Multivariate Time Series Data in Microservice Systems}

Metric-based anomaly detection and RCA are typically based on runtime information collected on the services within the microservice system. Such information includes metrics monitored on the microservices, such as workload, resource consumption, and response time. These metrics are typically represented as multivariate time series, with each time series corresponding to the data collected with a specific metric \cite{Soldani2022rcasurvey}. Microservices also generate logs to provide more detailed and meaningful information about their state. Some previous studies~\cite{wang2021detecting, aggarwal2020localization, aggarwal2021causal} parse raw logs to extract log static structures (i.e., log templates \cite{le2023log}), and count the occurrences of these templates, which are subsequently transformed into time series data. These logs can also be a source of multivariate time series data, which can be used for analyze root cause. In this work, we, however, only focus on metrics as multivariate time series data, as with~\cite{Li2022Circa, Azam2022rcd, wu2020microrca, wu2021microdiag, thalheim2017sieve, Meng2020Microcause, Wang2018cloudranger, liu2021microhecl}.


\vspace{-0.1cm}
\subsection{Anomaly Detection} \label{sec:background-ad}

In microservice systems, once a failure occurs in a service, it is typically reflected in the metrics data of that particular service, resulting in an anomaly or a change in its data distribution. Furthermore, a failure in one service can propagate across the systems and impact other services, subsequently leading to changes to the metrics data of those services as well. To detect a failure in a microservice system, the goal is to detect any anomalies or changes in the given metrics dataset \cite{Soldani2022rcasurvey}.


A wide range of anomaly detection methods exist for time series data \cite{BlazquezGarc2021ADreview}. In this paper, since we target the problem of identifying root causes for microservice systems, we only focus on root cause localization-oriented anomaly detectors employed or discussed in existing metric-based RCA studies \cite{lee2023eadro}. N-Sigma \cite{Jinjin2018Microscope} is used and discussed in MicroRank \cite{yu2021microrank}, CIRCA \cite{Li2022Circa}, Eadro \cite{lee2023eadro}, and MicroScope \cite{Jinjin2018Microscope}. SPOT \cite{siffer2017anomaly} is mentioned and evaluated in CIRCA \cite{Li2022Circa} and Eadro \cite{lee2023eadro}.  BIRCH \cite{zhang1996birch} is employed in MicroRCA \cite{wu2020microrca}  and MicroDiag \cite{wu2021microdiag}. Finally, Univariate Offline Bayesian Change Point Detection \cite{adams2007bayesian} is used in CauseInfer \cite{Chen2014Causeinfer, chen2016causeinfer}. We describe these methods in the below.

\paragraph{N-Sigma}
N-Sigma (i.e., the three-sigma rule of thumb) represents one of the simplest anomaly detection methods. It operates based on the assumption that the data points falling within three standard deviations of the mean of the data distribution are considered normal. Consequently, any data point $x$ that falls outside this range, i.e., $\mu - 3\sigma < x < \mu + 3\sigma$, is deemed abnormal. For instance, MicroScope \cite{Jinjin2018Microscope} computes the mean and standard deviation of the metrics data distribution using the most recent 10 minutes of data and then applies this method to detect anomalies.

\vspace{-0.2cm}
\paragraph{SPOT}
SPOT and dSPOT \cite{siffer2017anomaly} are founded based on the principles of Extreme Value Theory \cite{Coles2001EVT} to detect anomalies and have been employed in various RCA research works \cite{pan2021dycauserca, Li2022Circa, lee2023eadro, Meng2020Microcause}. SPOT is designed to handle data with stationary distribution, while dSPOT is developed for streaming data susceptible to concept drift. It is worth noting that different studies have used different variants of SPOT; For instance, MicroCause \cite{Meng2020Microcause} and DycauseRCA \cite{pan2021dycauserca} made use of dSPOT, whereas CIRCA~\cite{Li2022Circa} employed biSPOT. In this work, we will employ dSPOT, as in \cite{Meng2020Microcause, pan2021dycauserca}. Therefore, in the sequel, when referring to SPOT, we are specifically referring to dSPOT. 

\vspace{-0.2cm}
\paragraph{BIRCH}
BIRCH (Balanced Iterative Reducing and Clustering using Hierarchies) \cite{zhang1996birch} is a well-regarded unsupervised clustering algorithm known for its efficiency in real-time data analysis and anomaly detection, particularly in large-scale datasets and time-series data. Its main idea is to generate a brief and informative summary about the original dataset, and then perform clustering on this summary. BIRCH considers a data point to be an anomaly when it is in a different cluster with other consecutive data points. MicroRCA \cite{wu2020microrca} and MicroDiag \cite{wu2021microdiag} have employed BIRCH for anomaly detection as a preliminary step before conducting root cause analysis.

\vspace{-0.2cm}
\paragraph{Univariate Offline Bayesian Change Point Detection}
Bayesian change point detection \cite{adams2007bayesian} is a statistical method for identifying change points in time series data, i.e., timesteps where the data distribution experiences significant shifts. It relies on the principle of causal predictive filtering, aiming to generate an accurate distribution of future unseen data points based solely past observations. Through Bayesian statistics, it incorporates the prior information regarding the characteristics of the change points into the modelling process, making it to be both effective and efficient. In previous RCA research works, such as \cite{Chen2014Causeinfer} and \cite{chen2016causeinfer}, univariate offline Bayesian change point detection was used to detect change points (anomalies) within the time series metrics data. 

Besides the methods mentioned, commercial platforms like Datadog \cite{datadog} and Dynatrace \cite{dynatrace} also provide anomaly detection techniques \cite{datadog2, dynatrace1} for univariate time series. These approaches either rely on users or historical data to obtain thresholds to detect anomalies. They are similar to the concept of N-Sigma, which uses expected values along with pre-defined tolerance thresholds.

% Besides the methods mentioned, commercial platforms like Datadog \cite{datadog} and Dynatrace \cite{dynatrace} also provide threshold-based anomaly detection techniques \cite{datadog2, dynatrace1} for univariate time series. These approaches rely on pre-defined thresholds to detect anomalies in incoming data, similar to the concept of N-Sigma, which uses expected values along with pre-defined tolerance thresholds.


\subsection{Root Cause Analysis}

Existing metric-based RCA algorithms can be classified into three main categories: statistical analysis, topology graph-based methods, and causal graph-based methods \cite{Soldani2022rcasurvey}.

\vspace{-0.2cm}
\paragraph{Statistical Analysis.} These approaches pinpoint failures' root causes by identifying metrics that undergo significant changes during the anomalous period. $\epsilon$-Diagnosis \cite{Shan2019Ediagnosis} uses the two-sample test algorithm and $\epsilon$-statistics to measure the similarity among the metrics and rank the root cause based on the similarity scores. $\epsilon$-Diagnosis is evaluated against three statistical analysis methods: Pearson distance, KNN \cite{friedman1983graph, luo2014correlating} and MST \cite{friedman1979multivariate}. KNN \cite{friedman1983graph, luo2014correlating} uses nearest neighbours to model the distance between two time series, while MST \cite{friedman1979multivariate} uses a minimum spanning tree to represent the distance. In \cite{wang2020root}, a neural network is used to learn the normal behaviour and measure the similarity of monitoring data when the failure happens. They use mutual information to rank the root causes. N-Sigma \cite{Jinjin2018Microscope,Li2022Circa} is another statistical analysis technique that assesses the distance using z-score.

\vspace{-0.2cm}
\paragraph{Topology Graph-based Analysis.} These approaches reconstruct a topology graph representing the microservice system using information from monitoring systems and operational knowledge. MicroRCA \cite{wu2020microrca} constructs a topology graph from monitoring data, extracts anomalous subgraphs, and uses the random walk algorithm to infer root causes. Similarly, \cite{wu2020performance} constructs a topology graph and anomalous subgraphs, followed by a neural network-based method to infer the root cause. Sieve \cite{thalheim2017sieve} uses a clustering technique to reduce the number of metrics on the constructed topology graph, then uses Granger Causality tests \cite{Granger1980causal} to determine the possible root causes. Likewise, Brandon \cite{brandon2020graph} augments the provided topology graph with monitored metrics, conducts root cause searches through extracted subgraphs, and ranks root causes based on similarity scores. DLA \cite{samir2019dla} transforms the provided topology graph and metric data into a hierarchical hidden Markov model and identifies root causes by computing the path with the highest anomalous probability. Meanwhile, CIRCA \cite{Li2022Circa} performs hypothesis testing on the structural graph to find the root causes. Commercial platforms such as Datadog \cite{datadog} and Dynatrace \cite{dynatrace} construct a topology graph from distributed traces and perform Depth First Search (DFS) to find the root cause services \cite{datadog1, dynatrace3}.

\vspace{-0.2cm}
\paragraph{Causal Graph-based Analysis.} 

Many recent RCA techniques adopt the causal graph-based approach \cite{Chen2014Causeinfer, chen2016causeinfer, Jinjin2018Microscope, Wang2018cloudranger, Chen2019airalert, Ma2019Msrank, Meng2020Microcause, Ma2020Automap, wu2021microdiag, Azam2022rcd, Xin2023CausalRCA}. The main idea is to construct a causal graph where vertices represent services or metrics of the microservices, and edges represent the cause-effect relationships between the services/metrics. These graphs are constructed using different causal discovery methods such as PC, FCI, LiNGAM, and GES \cite{Spirtes1993Causal, Spirtes1995fci, Chickering2002Ges, Shimizu2006Lingam, Jaber2020PsiFCI}. Assuming the root cause metric would affect other services' metrics, graph centrality algorithms like Breath First Search (BFS), random walk, or PageRank are used to rank the root causes. In addition, correlation analysis can be conducted to measure the correlation among metrics \cite{wu2021microdiag}, and the graph traversal process can consider these scores to identify the root cause. Recently, RCD \cite{Azam2022rcd} employs a divide-and-conquer strategy to split the input metrics into smaller chunks and constructs a causal graph for each chunk. It then employs $\Psi$-PC \cite{Jaber2020PsiFCI} to identify the root cause for each chunk, which is later combined to yield the final rootcause. CausalRCA \cite{Xin2023CausalRCA} introduces the use of a gradient-based causal discovery method, namely DAG-GNN, to uncover causal relationships among metrics.
